\documentclass[a4paper,10pt]{article}

\usepackage[utf8]{inputenc}
\usepackage[T1]{fontenc}
\usepackage{geometry}
\usepackage{tabularx}
\usepackage{booktabs}
\usepackage{listings}
\usepackage{xcolor}
\usepackage{amsmath}
\usepackage[hidelinks]{hyperref}

\geometry{
    a4paper,
    top=2cm,
    bottom=2cm,
    left=2cm,
    right=2cm
}

\newcommand{\hex}[1]{\texttt{0x#1}}
\newcommand{\code}[1]{\texttt{#1}}

\lstdefinestyle{cstyle}{
    language={[ANSI]C},
    basicstyle=\ttfamily\small,
    keywordstyle=\bfseries,
    commentstyle=\itshape,
    numbers=left,
    numberstyle=\tiny,
    numbersep=7pt,
    frame=single,
    breaklines=true,
    showstringspaces=false,
    tabsize=4
}

\title{\textbf{DEOS Ethernet Transport Adapter (TCP \& UDP) ICD}\\
\large Sürüm 1.1 -- Textual (TCP) \& ROS 2 Bridge (UDP) Arayüzleri}
\author{\textbf{Emir Arkalı} (\href{mailto:emirarkali07@gmail.com}{emirarkali07@gmail.com})\\DEOS -- Dokuz Eylül Otonom Sistemler}
\date{İlk Yayın: 28 Eylül 2026 \quad | \quad Son Revizyon: 13 Ekim 2026}

\begin{document}

\maketitle

\begin{center}
\small
\begin{tabularx}{\textwidth}{ll lX}
\toprule
\textbf{Doküman Tipi:} & Interface Control Document (ICD) & \textbf{Doküman Kodu:} & \code{DOC-ICD-TRP-ETH-006} \\
\textbf{Doküman Sürümü:} & Sürüm 1.1 & \textbf{Protokol Sürümü:} & v1.1 (\hex{11}) -- Medya Bağımsız \\
\textbf{Yazar:} & Emir Arkalı (\texttt{emirarkali07@gmail.com}) & \textbf{Son Değiştiren:} & Emir Arkalı \\
\textbf{İlk Yayın Tarihi:} & 28 Eylül 2026 & \textbf{Son Revizyon Tarihi:} & 13 Ekim 2026 \\
\textbf{Durum:} & Yürürlükte / Onaylandı & \textbf{Taşıma / Kapsam:} & TCP (Port 5000) \& UDP (Port 5001) Taşıma Katmanı \\
\bottomrule
\end{tabularx}
\end{center}

\subsection*{Doküman Revizyon Geçmişi}
\begin{center}
\small
\begin{tabularx}{\textwidth}{c c l X}
\toprule
\textbf{Revizyon} & \textbf{Tarih} & \textbf{Yazar / Değiştiren} & \textbf{Açıklama} \\
\midrule
v0.3 & 28 Eylül 2026 & Emir Arkalı & İlk taslak ve monolitik ICD v0.3 tanımı. \\
v1.0 & 04 Ekim 2026 & Emir Arkalı & TCP/Ethernet adaptörü bağımsız doküman oluşturuldu. \\
v1.1 & 13 Ekim 2026 & Emir Arkalı & Sürüm 1.1 standardizasyonu. Yerel ROS 2 UDP Bridge (Port 5001) taşıma katmanı entegre edildi; Textual ve ROS Bridge Node ID eşleşmeleri düzeltildi. \\
\bottomrule
\end{tabularx}
\end{center}
\vspace{0.8em}

\section{Giriş ve Kapsam}
Bu doküman, DEOS ağındaki Araç Ana Kontrolcüsü (\code{DEOS\_NODE\_MAIN\_STM32}, \hex{03}) üzerinde koşan Ethernet taşıma katmanının kurallarını ve çerçeve (framing) yapısını tanımlar. 

DEOS Ethernet donanım hattı, üzerindeki tek fiziksel PHY/MAC üzerinden iki farklı mantıksal soket arayüzü sunar:
\begin{enumerate}
    \item \textbf{Textual Konsol Arayüzü (TCP Port 5000):} Operatör bilgisayarı veya Raspberry Pi üzerindeki Python/CLI arayüzü ile bağlantı odaklı, sıralı ve garantili akış haberleşmesi.
    \item \textbf{ROS 2 Köprü Arayüzü (UDP Port 5001):} Gömülü Linux otonom sürüş bilgisayarı üzerindeki \code{deos\_ros\_bridge} paketi ile çift yönlü, düşük gecikmeli, datagram tabanlı haberleşme.
\end{enumerate}

Ethernet bağlantısında fiziksel katman (L2) donanımsal 32-bit CRC (FCS) doğrulaması yaptığı ve L4 taşıma protokolleri (TCP/UDP) kendi dahili checksum mekanizmalarını barındırdığı için, LoRa tarafında kullanılan yazılımsal CRC16 hesaplamalarına bu hatta \textbf{ihtiyaç duyulmaz}. Bununla birlikte, TCP akışındaki olası bayt kaymalarına karşı koruma sağlamak amacıyla TCP çerçevelerinin başında \code{0xAA 0x55} senkronizasyon baytları yer alır; mesaj tabanlı çalışan UDP tarafında ise ek bir magic bayta gerek duyulmaz.

\section{Textual TCP Konsol Arayüzü (Port 5000)}

\subsection{Bağlantı Özellikleri}
\begin{itemize}
    \item \textbf{Protokol:} TCP Stream (\code{SOCK\_STREAM})
    \item \textbf{Roller:} Zephyr (STM32) kartı \textbf{TCP Server} olarak dinlemede kalır. Textual arayüzü (Python) açıldığı anda \textbf{TCP Client} olarak karta bağlanır.
    \item \textbf{Varsayılan IP \& Port:} \code{192.168.1.10:5000}
    \item \textbf{Kaynak/Hedef Düğüm:} \code{DEOS\_NODE\_TEXTUAL} (\hex{01})
\end{itemize}

\subsection{TCP Çerçeve (Frame) Formatı}
TCP veri akışında çerçeve sınırlarını belirlemek ve akış kaymalarına karşı anında yeniden senkronize olabilmek (Stream Resynchronization) için her paketin başına 2 baytlık \code{0xAA 0x55} sihirli baytları (Magic Word) ve 1 baytlık \code{Length} bilgisi eklenir:

\begin{table}[h]
\centering
\begin{tabular}{|c|c|c|c|}
\hline
\textbf{Magic Word (2 Bayt)} & \textbf{Length (1 Bayt)} & \textbf{DEOS Common Header (8 Bayt)} & \textbf{DEOS Payload ($N$ Bayt)} \\
\hline
\hex{AA} \hex{55} & \hex{08} - \hex{44} & Priority, Class, Service, Dest, Src, Ver, Seq, Cmd & Komut Verisi (Örn: Ping ID) \\
\hline
\end{tabular}
\end{table}

\begin{itemize}
    \item \textbf{Magic Word (2 Byte):} Daima \hex{AA} \hex{55} değerindedir. Alıcı akıştan koparsa bir sonraki \hex{AA} \hex{55} desenini arayarak akışı anında kurtarır.
    \item \textbf{Length (1 Byte):} Kendisinden sonra gelen DEOS Header ve Payload alanlarının toplam bayt sayısıdır (\code{8 + payload\_len}).
    \item \textbf{DEOS Common Header (8 Byte):} Sırasıyla \code{priority}, \code{message\_class}, \code{service}, \code{destination}, \code{source}, \code{version}, \code{sequence}, \code{command}.
    \item \textbf{DEOS Payload ($N$ Byte):} Gönderilmeyen boş baytlar hatta basılmaz; yalnızca dolu olan $N$ bayt aktarılır.
\end{itemize}

\section{ROS 2 Köprüsü UDP Arayüzü (Port 5001)}

\subsection{Bağlantı Özellikleri}
\begin{itemize}
    \item \textbf{Protokol:} UDP Datagram (\code{SOCK\_DGRAM})
    \item \textbf{Roller:} Zephyr (STM32) kartı tüm arayüzlerde (\code{INADDR\_ANY}) \textbf{5001} portunu dinler. Otonom kontrol tarafındaki \code{deos\_ros\_bridge} C++ node'u bu porta datagram yollar.
    \item \textbf{Varsayılan IP \& Port:} \code{192.168.1.10:5001}
    \item \textbf{Mantıksal Düğüm:} \code{DEOS\_NODE\_ROS\_BRIDGE} (\hex{04})
\end{itemize}

\subsection{Dinamik İstemci Keşfi (Dynamic Client Discovery)}
ROS 2 köprüsü statik bir IP'ye zorlanmak zorunda değildir. STM32 Main Node (\code{eth\_ros\_bridge.cpp}):
\begin{enumerate}
    \item Port 5001 üzerinden gelen ilk geçerli pakette göndericinin IP adresini ve port numarasını (\code{m\_client\_addr}) hafızaya alır (\textit{Discovered}).
    \item Araç içindeki CAN düğümlerinden veya sensörlerden ROS 2 köprüsüne (\code{DEOS\_NODE\_ROS\_BRIDGE}) gitmesi gereken tüm telemetri mesajlarını \code{zsock\_sendto} ile bu keşfedilen istemci adresine basar.
    \item İstemci adresi veya portu değişirse, gelen yeni paketle birlikte hedef adres otomatik güncellenir.
\end{enumerate}

\subsection{UDP Datagram Paket Yapısı}
UDP datagramları sınırları belirli paketler (packet-oriented) halinde taşınır. Her UDP datagramının başında 1 baytlık uzunluk göstergesi yer alır:

\begin{table}[h]
\centering
\begin{tabular}{|c|c|c|}
\hline
\textbf{Length (1 Bayt)} & \textbf{DEOS Common Header (8 Bayt)} & \textbf{DEOS Payload ($N$ Bayt)} \\
\hline
$8 + N$ & 8 Bayt Sabit Başlık & Dinamik Veri Baytları \\
\hline
\end{tabular}
\end{table}

\begin{itemize}
    \item Toplam datagram boyutu: \code{Length + 1} bayttır.
    \item Alıcı taraf, gelen datagramın ilk baytını okur (\code{rx\_len = rx\_buf[0]}); eğer alınan toplam bayt sayısı \code{rx\_len + 1} değerine eşit ve \code{rx\_len} $\ge$ 8 ise paketi doğrudan DEOS çekirdeğine iletir.
\end{itemize}

\section{Örnek Mesaj Anatomisi (PING Mesajı)}

İstemciden (Textual veya ROS 2 Bridge) karta atılan bir PING paketinin bayt dizilimi:

\begin{table}[h]
\centering
\small
\begin{tabular}{|c|l|c|l|}
\hline
\textbf{Byte Indeksi} & \textbf{Alan} & \textbf{Örnek (Hex)} & \textbf{Açıklama} \\
\hline
\textbf{0}  & Length & \hex{0C} & 12 Byte (8 Bayt Header + 4 Bayt Payload) \\
\textbf{1}  & priority & \hex{01} & \code{DEOS\_PRIO\_NETWORK} \\
\textbf{2}  & message\_class & \hex{01} & \code{DEOS\_CLASS\_NETWORK} \\
\textbf{3}  & service & \hex{00} & \code{DEOS\_SERVICE\_SYSTEM} \\
\textbf{4}  & destination & \hex{03} & \code{DEOS\_NODE\_MAIN\_STM32} \\
\textbf{5}  & source & \hex{04} & \code{DEOS\_NODE\_ROS\_BRIDGE} (\hex{01} Textual için) \\
\textbf{6}  & version & \hex{11} & Protokol Versiyonu (v1.1) \\
\textbf{7}  & sequence & \hex{05} & Mesaj Sıra Numarası \\
\textbf{8}  & command & \hex{01} & \code{DEOS\_CMD\_SYSTEM\_PING} \\
\textbf{9}  & payload[0] & \hex{45} & PING ID: 69 (Little Endian) LSB \\
\textbf{10} & payload[1] & \hex{00} & PING ID \\
\textbf{11} & payload[2] & \hex{00} & PING ID \\
\textbf{12} & payload[3] & \hex{00} & PING ID MSB \\
\hline
\end{tabular}
\end{table}

\section{Referans İstemci Kodları}

\subsection{Python TCP İstemci (Textual - Port 5000)}
\begin{lstlisting}[language=Python, basicstyle=\ttfamily\small, frame=single]
import socket, struct

sock = socket.socket(socket.AF_INET, socket.SOCK_STREAM)
sock.connect(("192.168.1.10", 5000))

# 8 bayt DEOS Header: Prio, Class, Srv, Dest(0x03), Src(0x01), Ver(0x11), Seq, Cmd(0x01)
header = struct.pack('<8B', 1, 1, 0, 0x03, 0x01, 0x11, 1, 0x01)
payload = struct.pack('<I', 69) # Ping ID: 69
msg = header + payload

magic = bytes([0xAA, 0x55])
length_byte = bytes([len(msg)]) # 12 Byte
sock.sendall(magic + length_byte + msg)
\end{lstlisting}

\subsection{C++ ROS 2 Bridge UDP İstemci (Port 5001)}
\begin{lstlisting}[style=cstyle]
// deos_ros_bridge icerisindeki udp_transport ornegi
uint8_t tx_buf[64];
uint8_t payload_len = 4;
uint8_t len = 8 + payload_len; // 12

tx_buf[0] = len;
tx_buf[1] = DEOS_PRIO_COMMAND;
tx_buf[2] = DEOS_CLASS_COMMAND;
tx_buf[3] = DEOS_SERVICE_TRACTION;
tx_buf[4] = DEOS_NODE_TRACTION;     // Hedef: 0x10
tx_buf[5] = DEOS_NODE_ROS_BRIDGE;   // Kaynak: 0x04
tx_buf[6] = 0x11;                   // DEOS v1.1
tx_buf[7] = seq++;
tx_buf[8] = 0x01;                   // SET_SPEED

// Payload (Hedef hiz int32: 500 -> 5.00 m/s)
int32_t target_speed = 500;
memcpy(&tx_buf[9], &target_speed, sizeof(target_speed));

sendto(udp_socket, tx_buf, len + 1, 0, 
       (struct sockaddr*)&stm32_addr, sizeof(stm32_addr));
\end{lstlisting}

\end{document}
